\clearpage\begingroup
\emergencystretch=4em\relax\NoSourceHead\SourceLeaf{480}
\makeatletter\tagsleft@true\let\veqno\@@leqno\makeatother
\renewcommand{\thefootnote}{\arabic{footnote})}\setcounter{footnote}{34}
\ArithmeticHead{23}{R}{GAUSS’ PRINCIPIA GENERALIA THEORIAE FIGURAE FLUIDORUM.}
\begin{center}\textls[100]{A) Aufstellung der ersten Variation.}\end{center}
\noindent\hspace*{\parindent}11. Für die Variationsrechnung weitaus am wichtigsten ist derjenige Teil
der \textsc{Gauss}schen Untersuchung, der sich mit der Berechnung der Variation des
Flächeninhalts $U$ beschäftigt.

Die in der Form
\begin{equation*}
z=z(x,y)\tag{28}
\end{equation*}
angenommene freie Oberfläche $U$ wird in der Weise variiert, dass jeder Punkt
$x,y,z$ der Fläche durch einen benachbarten Punkt\footnote{Die Bezeichnungen $\bar x$, $\bar y$, $\bar z$, $\bar U$, $\mathfrak{A}$ finden sich bei \textsc{Gauss} nicht; ebensowenig weiter unten die Bezeichnungen $M$, $M_1$, $M_2$, u. s. w.}
\begin{equation*}
\bar x=x+\delta x,\quad\bar y=y+\delta y,\quad\bar z=z+\delta z\tag{29}
\end{equation*}
ersetzt wird, wobei $\delta x$, $\delta y$, $\delta z$ »unbestimmte, aber unendlich kleine Funktionen
von $x$ und $y$« sind, deren erste partielle Ableitungen stillschweigend ebenfalls
als unendlich klein vorausgesetzt werden.

Die so entstehende variierte Fläche, ebenso wie ihr Inhalt, möge mit
$\bar U$ bezeichnet werden. Beide Flächenstücke $U$ und $\bar U$ werden erhalten, wenn
die unabhängigen Variabeln $x,y$ ein und denselben Bereich der $x,y$-Ebene
überstreichen, nämlich die Projektion $\mathfrak{A}$ des Flächenstücks $U$ auf die $x,y$-Ebene.

Die Variation, — worunter immer die erste Variation verstanden wird —,
$\delta U$ ist dann der Zuwachs $\bar U-U$ bei Vernachlässigung unendlich kleiner Grössen
höherer Ordnung. Bei der Berechnung von $\delta U$ geht nun \textsc{Gauss} in dem Bestreben,
möglichst wenig bei seiner Ableitung vorauszusetzen und den Beweis
von den allereinfachsten Elementen aus aufzubauen, so weit, dass er nicht
einmal den allgemeinen Ausdruck für den Inhalt eines Flächenstücks als bekannt
voraussetzt. Vielmehr berechnet er direkt den Inhalt eines Elementes
der ursprünglichen Fläche $U$ und den des entsprechenden Elementes der
variierten Fläche $\bar U$. Daraus erhält er zunächst die Variation des Flächenelementes
und erst weiterhin durch Integration die Variation des Gesamtflächeninhaltes
$U$. Im einzelnen verfährt \textsc{Gauss} folgendermassen (Art. 21):
\par\clearpage\SourceLeaf{481}
\ArithmeticHead{24}{L}{BOLZA, GAUSS UND DIE VARIATIONSRECHNUNG.}
\noindent\hspace*{\parindent}Sind $M$, $M_1$, $M_2$ die Punkte der Fläche $U$ bezw. mit den Abszissen
\[
(x,y),\quad (x+dx,\ y+dy),\quad (x+d'x,\ y+d'y)
\]
und $N_1$, $N_2$ diejenigen Punkte des Raumes, die aus $M_1$, $M_2$ hervorgehen, wenn
man in deren $z$-Koordinaten die vollständigen Inkremente von $z$ durch die
entsprechenden Differentiale ersetzt, also
\[
z(x+dx,\ y+dy)\quad\text{durch}\quad z(x,y)+p\,dx+q\,dy,
\]
bezw.
\[
z(x+d'x,\ y+d'y)\quad\text{durch}\quad z(x,y)+p\,d'x+q\,d'y,
\]
so nimmt \textsc{Gauss} das Dreieck $MN_1N_2$ als Element der ursprünglichen Fläche
und berechnet dessen Inhalt $dU$ nach der bekannten Formel für den Flächeninhalt
eines Dreiecks. Den drei Punkten $M$, $M_1$, $M_2$ der Fläche $U$ entsprechen
nun mittels der Zuordnung (29) auf der variierten Fläche $\bar U$ drei Punkte
$\bar M$, $\bar M_1$, $\bar M_2$; mit $\bar N_1$, $\bar N_2$ mögen diejenigen Punkte bezeichnet werden, die
aus $\bar M_1$, $\bar M_2$ hervorgehen, wenn man in den Ausdrücken der Koordinaten
wieder die vollständigen Inkremente von $z$, $\delta x$, $\delta y$, $\delta z$ durch die Differentiale
ersetzt. Dann betrachtet \textsc{Gauss} als Element der variierten Fläche das Dreieck
$\bar M\bar N_1\bar N_2$ und berechnet dessen Inhalt wieder nach der allgemeinen Formel
für den Inhalt eines Dreiecks. Nunmehr werden die höheren Potenzen von
$\delta x$, $\delta y$, $\delta z$ und ihren partiellen Ableitungen vernachlässigt; das so erhaltene
Resultat möge mit $d\bar U$ bezeichnet werden. Dann ist
\[
d\bar U-dU=\delta dU.
\]
Daraus ergibt sich schliesslich $\delta U$ durch Integration über die Fläche $U$.

Dabei führt \textsc{Gauss} an Stelle der partiellen Ableitungen $p,q$ die Richtungskosinus
$\xi$, $\eta$, $\zeta$ der in Beziehung auf die Flüssigkeit äusseren Normale der
Fläche $U$ im Punkt $x,y,z$ ein und erhält so schliesslich das Resultat in der
Form
\begin{equation*}
\begin{aligned}
\delta U={}&\int dU\left((\eta^2+\zeta^2)\frac{\partial\delta x}{\partial x}-\xi\eta\frac{\partial\delta y}{\partial x}-\xi\zeta\frac{\partial\delta z}{\partial x}\right)\\
&+\int dU\left(-\xi\eta\frac{\partial\delta x}{\partial y}+(\xi^2+\zeta^2)\frac{\partial\delta y}{\partial y}-\eta\zeta\frac{\partial\delta z}{\partial y}\right).
\end{aligned}\tag{30}
\end{equation*}

Wir bemerken, dass sich dasselbe Resultat ergibt, wenn man in der \textsc{Lagrange}schen
Fundamentalformel (21) $V=\sqrt{1+p^2+q^2}$ setzt.
\par\clearpage\SourceLeaf{482}
\ArithmeticHead{25}{R}{GAUSS’ PRINCIPIA GENERALIA THEORIAE FIGURAE FLUIDORUM.}
\noindent\hspace*{\parindent}12. Um zu erkennen, was das eben geschilderte \textsc{Gauss}sche Verfahren,
bei dem die Schwierigkeiten der Berechnung des Inhalts einer krummen Fläche
in eigentümlicher Weise mit denjenigen der Variation der Doppelintegrale
verwoben sind, für die \textls[100]{allgemeine} Variationsrechnung bedeutet, hat man
zunächst die beiden Arten von Schwierigkeiten von einander zu trennen, indem
man die allgemeine Formel für den Inhalt einer in Parameterdarstellung gegebenen
Fläche, die \textsc{Gauss}\footnote{\textit{Theoria attractionis corporum sphaeroidicorum etc.}, Comm. Soc. Reg. sc. Gotting. rec. vol. II, 1813, Art. 10; Werke, Bd. V, S. 15.} schon im Jahr 1813 angegeben hatte, als bekannt
und für die beiden Flächenstücke $U$ und $\bar U$ als gültig voraussetzt. Berechnet
man nach dieser Formel den Inhalt einerseits der durch die Gleichung
(28) dargestellten Fläche $U$, andererseits der durch die Gleichungen (29) in
Parameterdarstellung mit $x,y$ als Parametern dargestellten Fläche $\bar U$, wobei
in beiden Integralen die unabhängigen Variabeln auf denselben Bereich $\mathfrak{A}$ zu
beschränken sind, und bildet die Differenz $\bar U-U$, so erhält man nach Vernachlässigung
der Glieder höherer Ordnung in $\delta x$, $\delta y$, $\delta z$ und ihren partiellen
Ableitungen den \textsc{Gauss}schen Ausdruck (30) für $\delta U$. Das ist sozusagen der
der Variationsrechnung angehörige Bestandteil des \textsc{Gauss}schen Gedankengangs.

Wir versuchen jetzt, denselben Gedankengang auf das allgemeine Integral
\[
J=\iint\limits_{\mathfrak{A}}V(x,y,z,p,q)\,dx\,dy
\]
mit variablem Integrationsbereich zu übertragen. Wir haben dann zunächst
dasselbe Integral $J$ für die Fläche $\bar U$ zu bilden, d. h. wenn wir alle auf die
Fläche $\bar U$ bezüglichen Grössen durch Überstreichen auszeichnen,
\[
\bar J=\iint\limits_{\bar{\mathfrak{A}}}V(\bar x,\bar y,\bar z,\bar p,\bar q)\,d\bar x\,d\bar y.
\]
Dabei ist die Fläche $\bar U$ in der aus den Gleichungen (29) durch Elimination
von $x,y$ sich ergebenden Form
\begin{equation*}
\bar z=\bar z(\bar x,\bar y)\tag{31}
\end{equation*}
gedacht, $\bar p$, $\bar q$ sind die Ableitungen von $\bar z$ nach $\bar x$ und $\bar y$, und $\bar{\mathfrak{A}}$ ist die Projektion
von $\bar U$ auf die $x,y$-Ebene.

Kehrt man jetzt von der Darstellung (31) der Fläche $\bar U$ zur Parameterdarstellung
\par\smallskip\noindent{\small X 2 Abh. 5.}\hfill{\small 4}\par
\clearpage\SourceLeaf{483}
\ArithmeticHead{26}{L}{BOLZA, GAUSS UND DIE VARIATIONSRECHNUNG.}
\noindent (29) zurück, so bedeutet das, dass man in dem Doppelintegral $\bar J$
durch die Substitution
\[
\bar x=x+\delta x,\quad\bar y=y+\delta y
\]
$x,y$ als unabhängige Variable einführt. Das gibt aber
\[
\bar J=\iint\limits_{\mathfrak{A}}V(\bar x,\bar y,\bar z,\bar p,\bar q)
\left(\frac{\partial\bar x}{\partial x}\frac{\partial\bar y}{\partial y}-\frac{\partial\bar x}{\partial y}\frac{\partial\bar y}{\partial x}\right)dx\,dy,
\]
wobei nunmehr der Integrationsbereich für $\bar J$ derselbe ist wie für $J$.

Durch die Zurückführung des Integrationsbereichs des Integrals $\bar J$ auf
den Integrationsbereich von $J$ sind nun aber die mit der Variation des Integrationsbereichs
verbundenen grundsätzlichen Schwierigkeiten überwunden, und
die Berechnung von $\delta J$ bietet keine Schwierigkeiten mehr. Das Resultat ist
die \textsc{Lagrange}sche Formel (21) mit den früher angegebenen Werten (18), (16)
für $\delta V$, $\delta p$, $\delta q$. Diese Verallgemeinerung des \textsc{Gauss}schen Verfahrens ist zwar
im Grunde mit der \textsc{Poisson}schen Methode der Parameterdarstellung aufs engste
verwandt, verdient aber doch als besondere Methode der Behandlung der
Doppelintegrale aufgeführt zu werden; wir werden sie die \textls[100]{Substitutionsmethode}
nennen. Sie findet sich zuerst klar ausgesprochen und durchgeführt
in der (trotz einiger Fehler) vortrefflichen Darstellung der Variationsrechnung,
die \textsc{Bordoni}\footnote{\textsc{A. M. Bordoni}, \textit{Lezioni di calcolo sublime}, Mailand 1831, Bd. II, Art. 398.} im Anschluss an die Darstellung von \textsc{Lagrange}\begingroup\renewcommand{\thefootnote}{37a)}\footnote{\textsc{J. L. Lagrange}, \textit{Leçons sur le calcul des fonctions}, Journal de l'École Polyt., cah. 12, 1804, nouvelle édition 1806, Oeuvres t. X.}\endgroup\setcounter{footnote}{37}\ und diese
nach den verschiedenen Richtungen weiterführend gegeben hat, und die in
der Literatur der Variationsrechnung viel zu wenig Beachtung gefunden hat.

Die Entwicklungen von \textsc{Bordoni} sind allem Anschein nach unbeeinflusst
durch \textsc{Gauss} und vermutlich ohne Kenntnis der \textsc{Gauss}schen Arbeit geschrieben.

Später ist dieselbe Methode, verallgemeinert, von \textsc{Sabinine}\footnote{\textit{Démonstration d'une formule de M. Ostrogradsky relative au calcul des variations des intégrales multiples}, Journal für r. u. a. Mathematik, Bd. 59 (1861), S. 185—189.} zum Beweis
der beim $n$-fachen Integral der \textsc{Lagrange}schen Fundamentalformel entsprechenden
Formel von \textsc{Ostrogradsky} und \textsc{Pagani} benutzt worden. Auch
\textsc{C. Jordan} legt diese Methode seiner Darstellung der Variation mehrfacher
Integrale\footnote{\textsc{C. Jordan}, \textit{Cours d'Analyse de l'École Polyt.}, Bd. III (1896), Nr. 396.} zugrunde.
\par\clearpage\endgroup
